%! Unit 3: Polynomial Functions — Summative Assessment — Practice Test --- Study Copy (blank)
\documentclass[10pt]{article}
\usepackage{algebra2-article}
\usepackage{algebra2-boxes}
\usepackage{grading}
\gradingassignment{alg2-u03-practice-test}
\gradingcourse{algebra_2}
\usepackage{algebra2-graded}
% Coordinate grid: {xmin}{xmax}{ymin}{ymax}
\newcommand{\lgrid}[4]{%
  \draw[step=1, linegray!40, very thin] (#1,#3) grid (#2,#4);
  \draw[->, semithick, charcoal] (#1,0)--(#2,0) node[below right, font=\scriptsize]{$x$};
  \draw[->, semithick, charcoal] (0,#3)--(0,#4) node[left, font=\scriptsize]{$y$};}
\newcommand{\dpt}[2]{\fill[charcoal] (#1,#2) circle (3pt);}

% Part-divider strip (headlinebox takes one arg: the background color)
\newcommand{\parthead}[1]{%
  \vspace{6pt}\noindent
  \begin{headlinebox}{forest}\color{white}\bfseries #1\end{headlinebox}%
  \vspace{4pt}}

\begin{document}

\pageheader{Unit 3: Polynomial Functions}{Practice Test --- Study Copy}
\namedateperiod

\vspace{0.10in}
\begin{remindbox}
\small
\textbf{This is a practice test.} It mirrors the real Unit 3 test in length, format,
and the ideas it covers, but the numbers and contexts differ. Work every problem
with no notes, then check your answers against the key.
\end{remindbox}

% ======================================================================
\parthead{Part A --- Vocabulary (8 pts)}
Write the letter of the best description next to each term.

\vspace{4pt}
\noindent
\begin{minipage}[t]{0.40\textwidth}
\begin{enumerate}[leftmargin=2.2em, itemsep=7pt, label=\arabic*.]
  \item Degree \inblank{\blank{1.2cm}}{\gradedblank[id=qa1, type=short, points=1, answer={B}, check=text]{1.2cm-2pt}{10pt}}
  \item Leading coefficient \inblank{\blank{1.2cm}}{\gradedblank[id=qa2, type=short, points=1, answer={A}, check=text]{1.2cm-2pt}{10pt}}
  \item End behavior \inblank{\blank{1.2cm}}{\gradedblank[id=qa3, type=short, points=1, answer={E}, check=text]{1.2cm-2pt}{10pt}}
  \item Multiplicity of a zero \inblank{\blank{1.2cm}}{\gradedblank[id=qa4, type=short, points=1, answer={C}, check=text]{1.2cm-2pt}{10pt}}
  \item Turning point \inblank{\blank{1.2cm}}{\gradedblank[id=qa5, type=short, points=1, answer={D}, check=text]{1.2cm-2pt}{10pt}}
  \item Factor Theorem \inblank{\blank{1.2cm}}{\gradedblank[id=qa6, type=short, points=1, answer={F}, check=text]{1.2cm-2pt}{10pt}}
  \item Rational Root Theorem \inblank{\blank{1.2cm}}{\gradedblank[id=qa7, type=short, points=1, answer={G}, check=text]{1.2cm-2pt}{10pt}}
  \item Fundamental Theorem of Algebra \inblank{\blank{1.2cm}}{\gradedblank[id=qa8, type=short, points=1, answer={H}, check=text]{1.2cm-2pt}{10pt}}
\end{enumerate}
\end{minipage}%
\hfill
\begin{minipage}[t]{0.57\textwidth}
\small
\begin{description}[leftmargin=1.4em, itemsep=3pt]
  \item[(A)] The number multiplying the highest-power term; its sign decides the right-hand arm.
  \item[(B)] The greatest exponent on the variable once the polynomial is in standard form.
  \item[(C)] How many times the factor $(x-r)$ appears --- odd means the graph crosses at $r$, even means it touches.
  \item[(D)] A point where the graph changes from increasing to decreasing, or the reverse.
  \item[(E)] What the graph does at its far left and far right, as $x\to\pm\infty$.
  \item[(F)] $(x-r)$ is a factor of $f(x)$ exactly when $f(r)=0$.
  \item[(G)] Every rational zero of a polynomial with integer coefficients looks like $\pm\frac{p}{q}$, where $p$ divides the constant term and $q$ divides the leading coefficient.
  \item[(H)] A polynomial of degree $n\ge 1$ has exactly $n$ complex zeros, counted with multiplicity.
\end{description}
\end{minipage}

% ======================================================================
\parthead{Part B --- Multiple Choice (12 pts)}
Circle the best answer.

\begin{enumerate}[leftmargin=1.9em, itemsep=9pt]
  \item \leavevmode\llap{\raisebox{0pt}[0pt][0pt]{\gradedblank[id=qb1, type=choice, points=2, answer={B}]{0.4in}{10pt}}\hspace{\dimexpr\leftmargin+4pt\relax}}The end behavior of $f(x)=-2x^4+3x^2-1$ is
    \hfill (A) both arms up \quad (B) both arms down \quad (C) left down, right up \quad (D) left up, right down
  \item \leavevmode\llap{\raisebox{0pt}[0pt][0pt]{\gradedblank[id=qb2, type=choice, points=2, answer={C}]{0.4in}{10pt}}\hspace{\dimexpr\leftmargin+4pt\relax}}$(2x-5)^2=$
    \hfill (A) $4x^2-25$ \quad (B) $4x^2+25$ \quad (C) $4x^2-20x+25$ \quad (D) $4x^2-10x+25$
  \item \leavevmode\llap{\raisebox{0pt}[0pt][0pt]{\gradedblank[id=qb3, type=choice, points=2, answer={A}]{0.4in}{10pt}}\hspace{\dimexpr\leftmargin+4pt\relax}}Factored completely, $x^3+27$ is
    \hfill (A) $(x+3)(x^2-3x+9)$ \quad (B) $(x+3)(x^2+3x+9)$ \quad (C) $(x+3)^3$ \quad (D) prime
  \item \leavevmode\llap{\raisebox{0pt}[0pt][0pt]{\gradedblank[id=qb4, type=choice, points=2, answer={A}]{0.4in}{10pt}}\hspace{\dimexpr\leftmargin+4pt\relax}}The remainder when $f(x)=x^3-4x^2+1$ is divided by $x-2$ is
    \hfill (A) $-7$ \quad (B) $-23$ \quad (C) $7$ \quad (D) $0$
  \item \leavevmode\llap{\raisebox{0pt}[0pt][0pt]{\gradedblank[id=qb5, type=choice, points=2, answer={C}]{0.4in}{10pt}}\hspace{\dimexpr\leftmargin+4pt\relax}}Which value is \emph{not} a possible rational zero of $f(x)=3x^3+2x^2-x+4$?
    \hfill (A) $-2$ \quad (B) $\frac{4}{3}$ \quad (C) $\frac{3}{4}$ \quad (D) $1$
  \item \leavevmode\llap{\raisebox{0pt}[0pt][0pt]{\gradedblank[id=qb6, type=choice, points=2, answer={A}]{0.4in}{10pt}}\hspace{\dimexpr\leftmargin+4pt\relax}}A polynomial with \emph{real} coefficients has degree $4$ and zeros $2$, $-1$, and $3i$. Its fourth zero must be
    \hfill (A) $-3i$ \quad (B) $3i$ \quad (C) $-3$ \quad (D) $\frac13$
\end{enumerate}

% ======================================================================
\parthead{Part C --- Short Answer \& Computation (40 pts)}
Show your work.

\begin{enumerate}[leftmargin=1.9em, itemsep=10pt]
  \item \textbf{Read the graph.} The polynomial function $f$ is graphed below. Each gridline is one unit.
        \begin{center}
        \begin{tikzpicture}[scale=0.5]
          \lgrid{-3}{4}{-5}{5}
          \begin{scope}
            \clip (-3,-5) rectangle (4,5);
            \draw[very thick, forest, domain=-1.5:3.2, samples=180, smooth]
              plot (\x,{0 - (\x)*(\x)*(\x) + 3*(\x)*(\x) - 4});
          \end{scope}
          \foreach \x in {-2,-1,1,2,3}
            \draw[charcoal] (\x,-0.15)--(\x,0.15) node[below=1pt, font=\tiny]{$\x$};
          \dpt{-1}{0}
          \dpt{2}{0}
          \dpt{0}{-4}
        \end{tikzpicture}
        \end{center}
        (a) Degree: even or odd? \inblank{\blank{1.4cm}}{\gradedblank[id=qc1a, type=short, points=1, answer={odd}, check=text]{1.4cm-2pt}{10pt}} \quad Leading coefficient: positive or negative? \inblank{\blank{1.8cm}}{\gradedblank[id=qc1b, type=short, points=1, answer={negative}, check=text]{1.8cm-2pt}{10pt}}
        \\[4pt]
        (b) As $x\to-\infty$, $f(x)\to$ \inblank{\blank{1.4cm}}{\gradedblank[id=qc1c, type=short, points=1, answer={+inf}, check=text, display={$+\infty$}]{1.4cm-2pt}{10pt}}; \quad as $x\to+\infty$, $f(x)\to$ \inblank{\blank{1.4cm}}{\gradedblank[id=qc1d, type=short, points=1, answer={-inf}, check=text, display={$-\infty$}]{1.4cm-2pt}{10pt}}
        \\[4pt]
        (c) Zeros, each with its multiplicity, and whether the graph crosses or touches there:
        \inblank{\blank{7.0cm}}{\hbox to 7.0cm{\gradedblank[id=qc1e, type=short, points=2, answer={x=-1 (multiplicity 1, crosses); x=2 (multiplicity 2, touches)}, check=text, display={$x=-1$ (multiplicity 1, crosses); $x=2$ (multiplicity 2, touches)}]{4.2cm}{10pt}\hfill}}
        \\[4pt]
        (d) $y$-intercept: \inblank{\blank{1.8cm}}{\gradedblank[id=qc1f, type=short, points=1, answer={(0,-4)}, check=text]{1.8cm-2pt}{10pt}}; \quad number of turning points: \inblank{\blank{1.2cm}}{\gradedblank[id=qc1g, type=short, points=1, answer={2}, check=numeric]{1.2cm-2pt}{10pt}}
        \\[4pt]
        (e) Relative maximum: \inblank{\blank{2.0cm}}{\gradedblank[id=qc1h, type=short, points=1, answer={(2,0)}, check=text]{2.0cm-2pt}{10pt}}; \quad relative minimum: \inblank{\blank{2.0cm}}{\gradedblank[id=qc1i, type=short, points=1, answer={(0,-4)}, check=text]{2.0cm-2pt}{10pt}};
        \\[4pt]
        \phantom{(e)} absolute maximum or minimum: \inblank{\blank{4.0cm}}{\gradedblank[id=qc1j, type=short, points=1, answer={none}, check=text, display={none --- odd degree, so the arms run to $\pm\infty$}]{4.0cm-2pt}{10pt}}
        \\[4pt]
        (f) Interval(s) where $f$ is increasing: \inblank{\blank{2.6cm}}{\gradedblank[id=qc1k, type=short, points=1, answer={(0,2)}, check=text]{2.6cm-2pt}{10pt}}; \quad domain: \inblank{\blank{2.0cm}}{\gradedblank[id=qc1l, type=short, points=1, answer={all reals}, check=text]{2.0cm-2pt}{10pt}}; \quad range: \inblank{\blank{2.0cm}}{\gradedblank[id=qc1m, type=short, points=1, answer={all reals}, check=text]{2.0cm-2pt}{10pt}}

  \item \textbf{Operations.} Simplify. Write each answer in standard form.
        \\[4pt]
        (a) $(3x^3-5x+2)-(x^3+4x^2-5)$
        \vspace{1.5cm}\vadjust{\workanswer{\gradedblank[id=qc2a, type=short, points=2, answer={2x^3-4x^2-5x+7}]{1.3in}{10pt}}}
        \\
        (b) $(2x-3)(x^2+4x-1)$
        \vspace{1.6cm}\vadjust{\workanswer{\gradedblank[id=qc2b, type=short, points=2, answer={2x^3+5x^2-14x+3}]{1.3in}{10pt}}}
        \\
        (c) For your answer to (b), give the degree \inblank{\blank{1.0cm}}{\gradedblank[id=qc2c, type=short, points=1, answer={3}, check=numeric]{1.0cm-2pt}{10pt}}, the leading coefficient \inblank{\blank{1.0cm}}{\gradedblank[id=qc2d, type=short, points=1, answer={2}, check=numeric]{1.0cm-2pt}{10pt}},
        and the end behavior \inblank{\blank{4.2cm}}{\gradedblank[id=qc2e, type=short, points=1, answer={left down, right up}, check=text]{4.2cm-2pt}{10pt}}

  \item \textbf{Factor completely.} Every factor must be prime over the integers.
        \\[4pt]
        (a) $2x^3-18x$ \hspace{1.5cm} (b) $27x^3-8$ \hspace{1.5cm} (c) $x^3+3x^2-4x-12$
        \vspace{2.8cm}\vadjust{\workanswer[(a)]{\gradedblank[id=qc3a, type=short, points=1, answer={2x(x-3)(x+3)}]{1.3in}{10pt}\quad\textbf{(b)}\ \gradedblank[id=qc3b, type=short, points=1, answer={(3x-2)(9x^2+6x+4)}]{1.8in}{10pt}\quad\textbf{(c)}\ \gradedblank[id=qc3c, type=short, points=1, answer={(x+3)(x-2)(x+2)}]{1.6in}{10pt}}}

  \item \textbf{Division and the Remainder Theorem.}
        \\[4pt]
        \answerbox{\gradedblank[id=qc4a, type=short, points=2, answer={x^2+2x+6+3/(x-2)}, display={$x^2+2x+6+\frac{3}{x-2}$}]{1.3in}{10pt}}(a) Divide using \emph{long division} and write the answer as $Q(x)+\dfrac{R}{D(x)}$:
        \[ (x^3+2x-9)\div(x-2) \]
        \vspace{2.6cm}
        \\
        (b) Use the \emph{Remainder Theorem} to find the remainder when $g(x)=x^3+2x^2-5x+6$ is divided
        by $x+3$. Is $x+3$ a factor of $g$? Explain in one sentence.
        \vadjust{\workarea{\gradedblank[id=qc4b, type=rubric, points=2, rubric={../rubrics/alg2-u03-practice-test-qc4b.yaml}, display={$x+3=x-(-3)$, so the remainder is $g(-3)=-27+18+15+6=12$. Since the remainder is $12$ and not $0$, $x+3$ is \emph{not} a factor (the Factor Theorem requires $g(-3)=0$).}]{\linewidth}{2.0cm-6pt}}}\vspace{2.0cm}

  \item \textbf{Synthetic division.} You are told that $x=3$ is a zero of
        \[ f(x)=x^3-2x^2-5x+6. \]
        Use synthetic division to divide out that zero, factor $f$ completely, and list \emph{all} of
        its zeros.
        \vspace{3.0cm}\vadjust{\workanswer[Factored:]{\gradedblank[id=qc5a, type=short, points=1, answer={(x-3)(x+2)(x-1)}]{1.6in}{10pt}\quad\textbf{Zeros:}\ \gradedblank[id=qc5b, type=short, points=1, answer={3, -2, 1}, check=set, display={$3,\ -2,\ 1$}]{1.0in}{10pt}}}

  \item \textbf{Rational Root Theorem.} For $f(x)=2x^3-x^2-8x+4$:
        \\[4pt]
        (a) List every possible rational zero.
        \vspace{1.5cm}\vadjust{\workanswer{\gradedblank[id=qc6a, type=short, points=1, answer={1, -1, 2, -2, 4, -4, 1/2, -1/2}, check=set, display={$\pm1,\ \pm2,\ \pm4,\ \pm\tfrac12$}]{1.8in}{10pt}}}
        \\
        (b) Find all real zeros of $f$, showing the test and the division you used.
        \vspace{3.0cm}\vadjust{\workanswer[Zeros:]{\gradedblank[id=qc6b, type=short, points=2, answer={2, 1/2, -2}, check=set, display={$2,\ \tfrac12,\ -2$}]{1.3in}{10pt}}}

  \item \textbf{Solve over the complex numbers.} $x^3-3x^2+4x-12=0$
        \\[4pt]
        (a) Before solving: how many solutions must this equation have, counted with multiplicity, and
        which theorem says so? \inblank{\blank{6.0cm}}{\gradedblank[id=qc7a, type=rubric, points=2, rubric={../rubrics/alg2-u03-practice-test-qc7a.yaml}, display={Exactly $3$, by the Fundamental Theorem of Algebra}]{6.0cm-2pt}{10pt}}
        \\[4pt]
        (b) Solve. Give every solution.
        \vspace{2.6cm}\vadjust{\workanswer[Solutions:]{\gradedblank[id=qc7b, type=short, points=2, answer={3, 2i, -2i}, check=set, display={$3,\ 2i,\ -2i$}]{1.3in}{10pt}}}
        \\
        (c) How many of your solutions are real, and how many $x$-intercepts does the graph of
        $y=x^3-3x^2+4x-12$ therefore have? \inblank{\blank{5.0cm}}{\gradedblank[id=qc7c, type=short, points=1, answer={1 real solution; 1 x-intercept}, check=text, display={One real solution ($x=3$), so exactly one $x$-intercept --- a graph can never show the imaginary pair}]{5.0cm-2pt}{10pt}}

  \item \textbf{Build a polynomial from its zeros.} Write a polynomial function of \emph{least degree}
        with \emph{real} coefficients whose zeros are $-1$ and $2i$. Give your answer in factored form
        \emph{and} in standard form.
        \vspace{2.8cm}\vadjust{\workanswer[Factored:]{\gradedblank[id=qc8a, type=short, points=1, answer={(x+1)(x^2+4)}]{1.3in}{10pt}\quad\textbf{Standard:}\ \gradedblank[id=qc8b, type=short, points=1, answer={x^3+x^2+4x+4}]{1.5in}{10pt}}}
\end{enumerate}

% ======================================================================
\parthead{Part D --- Extended Response (12 pts)}
Explain your reasoning in full sentences.

\begin{enumerate}[leftmargin=1.9em, itemsep=16pt]
  \item \textbf{The graph plan.} Consider $f(x)=-2(x+2)(x-1)^2(x-3)$. Do \emph{not} expand it.
        \textbf{(a)} State the degree and the leading coefficient, and explain how you got each one
        from the factored form. What end behavior do they force?
        \textbf{(b)} List the zeros with their multiplicities, and say whether the graph crosses or
        touches the $x$-axis at each.
        \textbf{(c)} Find the $y$-intercept.
        \textbf{(d)} Build a sign chart and state every interval on which $f(x)>0$.
        \textbf{(e)} At most how many turning points can this graph have, and why?
        \vadjust{\workarea{\gradedblank[id=qd1, type=rubric, points=6, rubric={../rubrics/alg2-u03-practice-test-qd1.yaml}, display={(a) Degrees add: $1+2+1=4$, so the degree is $4$. Leading coefficients multiply:
        $-2\cdot 1\cdot 1\cdot 1=-2$. An even degree with a negative leading coefficient sends
        \emph{both} arms down: $f(x)\to-\infty$ at each end.\\[3pt]
        (b) $x=-2$ (multiplicity 1, crosses); $x=1$ (multiplicity 2, touches); $x=3$ (multiplicity 1,
        crosses).\\[3pt]
        (c) $f(0)=-2(2)(-1)^2(-3)=12$, so the $y$-intercept is $(0,12)$.\\[3pt]
        (d) Test one value in each interval, keeping only signs: $f(-3)<0$, $f(0)=12>0$, $f(2)=8>0$,
        $f(4)<0$. So $f(x)>0$ on $(-2,1)\cup(1,3)$ --- that is, for $-2<x<3$ except $x=1$ --- and
        $f(x)<0$ on $(-\infty,-2)\cup(3,\infty)$. The sign does \emph{not} flip at $x=1$; that is the
        fingerprint of the even multiplicity.\\[3pt]
        (e) At most $4-1=3$ turning points, since a degree-$n$ polynomial turns at most $n-1$ times.}]{\linewidth}{6.0cm-6pt}}}\vspace{6.0cm}

  \item \textbf{Modeling.} An open-top box is made from a $12$ in by $6$ in sheet of cardboard by
        cutting a square of side $x$ from each corner and folding up the sides.
        \textbf{(a)} Write the volume $V(x)$ in factored form, then expand it into standard form.
        \textbf{(b)} State the feasible domain of the model and justify it from the sheet's dimensions.
        \textbf{(c)} Find $V(1)$ and state its units.
        \textbf{(d)} Check that $x=6$ makes $V(x)=0$. Explain why $6$ is a genuine zero of the
        polynomial but not a possible cut length.
        \vadjust{\workarea{\gradedblank[id=qd2, type=rubric, points=6, rubric={../rubrics/alg2-u03-practice-test-qd2.yaml}, display={(a) Cutting $x$ from both ends of each side leaves a base of $(12-2x)$ by $(6-2x)$ and a
        height of $x$, so $V(x)=x(12-2x)(6-2x)$. Expanding:
        $(12-2x)(6-2x)=4x^2-36x+72$, so $V(x)=4x^3-36x^2+72x$.\\[3pt]
        (b) Every dimension must be positive: $x>0$, $12-2x>0$ (so $x<6$), and $6-2x>0$ (so $x<3$).
        The tightest of these wins, so the feasible domain is $0<x<3$ --- the short side is what
        limits the cut.\\[3pt]
        (c) $V(1)=1(10)(4)=40$ cubic inches (in$^3$). Check against standard form:
        $4-36+72=40$ \checkmark\\[3pt]
        (d) $V(6)=6(12-12)(6-12)=6\cdot 0\cdot(-6)=0$, so $6$ really is a zero of the polynomial. But a
        $6$-inch cut from a $6$-inch side removes the entire width, and $6-2(6)=-6$ is a negative
        length --- there is no box at all. The polynomial does not know about cardboard; only the
        feasible domain $0<x<3$ describes a real box.}]{\linewidth}{5.5cm-6pt}}}\vspace{5.5cm}
\end{enumerate}

\end{document}
